\section{Recovery transfer to laws and consequence}
\label{sec:recovery-consequence-transfer}

Gate~8 shows that the split CPEL translation represents the complete four-valued
evaluation and that the classical sector is its Boolean anti-diagonal. Gate~9 asks
what may be transferred from this local representation to familiar logical laws and
consequence. The key methodological point is to distinguish three quantifier levels:
\[
\begin{array}{c}
\text{one formula at one model/world point}\\[2pt]
\Downarrow\\[2pt]
\text{one formula throughout one fixed model}\\[2pt]
\Downarrow\\[2pt]
\text{consequence over every model and world.}
\end{array}
\]
An implication between the first two levels does not by itself establish the third,
and none of the results below is a completeness theorem for either the source or
target proof system.

\subsection{The two independent law boundaries}

Write \(v=\llbracket\varphi\rrbracket_{m,w}\). Tolerant excluded middle means that
\(\varphi\lor\neg\varphi\) has positive support. Universal local LP explosion means
that \(\varphi\land\neg\varphi\) LP-entails every possible FDE conclusion value at
that point. The latter quantifies over values rather than formulas so that the
boundary does not depend on whether a particular fixed language and valuation happen
to denote all four values.

\begin{theorem}[Exact local recovery split]
\label{thm:gate9-local-laws}
For every \(m,w,\varphi\),
\[
\boxed{
\begin{aligned}
\operatorname{TLEM}_{m,w}(\varphi)
&\Longleftrightarrow \operatorname{GapFree}(v),\\
\operatorname{ULPExp}_{m,w}(\varphi)
&\Longleftrightarrow \operatorname{GlutFree}(v),\\
\operatorname{Classical}(v)
&\Longleftrightarrow
\operatorname{TLEM}_{m,w}(\varphi)
\land
\operatorname{ULPExp}_{m,w}(\varphi).
\end{aligned}
}
\]
Moreover, strict excluded middle is exact full recovery:
\[
\llbracket\varphi\lor\neg\varphi\rrbracket_{m,w}=\Tval
\quad\Longleftrightarrow\quad
\operatorname{Classical}(v).
\]
\statusverified
\end{theorem}

Thus LEM and LP explosion do not return as one indivisible classical package.
Tolerant LEM removes \(\Nval\) but permits \(\Bval\); universal LP explosion removes
\(\Bval\) but permits \(\Nval\). Their conjunction removes both nonclassical
vertices, and strict-\(\Tval\) LEM already characterizes that conjunction.

Gate~8 makes the same result visible directly in the split coordinates
\[
p=\llbracket\operatorname{tr}_{+}(\varphi)\rrbracket^C_{m,w},
\qquad
n=\llbracket\operatorname{tr}_{-}(\varphi)\rrbracket^C_{m,w}.
\]

\begin{corollary}[Coordinate form]
\label{cor:gate9-coordinate-laws}
At every model/world point,
\[
\boxed{
\begin{aligned}
\operatorname{TLEM}_{m,w}(\varphi)
&\Longleftrightarrow (p=1\lor n=1),\\
\operatorname{ULPExp}_{m,w}(\varphi)
&\Longleftrightarrow \neg(p=1\land n=1),\\
\operatorname{StrictLEM}_{m,w}(\varphi)
&\Longleftrightarrow n=\neg p.
\end{aligned}
}
\]
\statusverified
\end{corollary}

Completeness of the evidence channels supplies tolerant LEM, consistency supplies LP
explosion, and the anti-diagonal supplies both. This is a law-level reading of the
same Boolean square, not a new target semantics.

\subsection{Transfer from probability recovery}

The pointwise equivalences lift by universal quantification over the worlds of one
fixed model. In particular, the Gate~7 recursive condition
\texttt{Formula.ProbabilityRecoveryAdmissible}, together with
\texttt{ModelProbabilityIntegrity}, already yields classicality of the chosen
formula at every world. Gate~9 composes this result with the exact law boundaries.

\begin{corollary}[Probability recovery restores both laws]
\label{cor:gate9-probability-laws}
If \(m\) satisfies probability integrity and \(\varphi\) is probability-recovery
admissible in \(m\), then throughout \(m\):
\[
\llbracket\varphi\lor\neg\varphi\rrbracket=\Tval
\]
and \(\varphi\land\neg\varphi\) LP-entails every FDE conclusion value.
\statusverified
\end{corollary}

This is a formula- and model-relative theorem. Its premises must be proved for the
formula in question; it does not state that every formula in every integrity model
is recovered. Lockean incompleteness at a belief node remains a possible obstruction
by Gate~7.

\subsection{When ST and LP consequence coincide}

Let \(\models^{m}_{\mathrm{ST}}\) and \(\models^{m}_{\mathrm{LP}}\) quantify over all
worlds of one fixed model. LP consequence always implies ST consequence, since
strict truth of the antecedent includes positive designation. The reverse direction
needs only classicality of the antecedent throughout that model: a positively
designated classical antecedent must be exactly \(\Tval\).

\begin{theorem}[Model-relative consequence collapse]
\label{thm:gate9-model-consequence-collapse}
If \(\varphi\) is classical at every world of \(m\), then for every \(\psi\),
\[
\boxed{
\varphi\models^{m}_{\mathrm{ST}}\psi
\quad\Longleftrightarrow\quad
\varphi\models^{m}_{\mathrm{LP}}\psi.
}
\]
No classicality premise on the consequent is required for this equivalence.
\statusverified
\end{theorem}

One can also quantify over every model/world point while restricting attention at
each point to evaluations where the displayed formulas are classical. On this
\emph{classical-restricted} domain, the ST and LP definitions again coincide. This
is a semantic restriction of their domains; it is not an assertion that unrestricted
ST and LP consequence coincide over all four-valued models.

\subsection{Explosion exposes the global boundary}

The consequence regime matters independently of recovery. A direct contradiction
\(v\land\neg v\) is never strictly \(\Tval\), for any of the four input values.
Consequently, direct contradiction entails every formula under unrestricted ST
consequence: its strict antecedent condition is never met.

LP consequence is different. At \(v=\Bval\), the direct contradiction is again
\(\Bval\), hence designated. A one-world model in the artifact assigns one atom
\(\Bval\) and a second atom \(\Fval\), verifying
\[
(p\land\neg p)\not\models_{\mathrm{LP}}q.
\]
At every point where the displayed values are recovered, the same explosion schema
is valid. The artifact packages both facts into one separation theorem.

\begin{theorem}[Recovered explosion is not global LP explosion]
\label{thm:gate9-explosion-boundary}
There are formulas for which classical-restricted LP explosion holds while
unrestricted LP consequence fails:
\[
\boxed{
(p\land\neg p)\models^{\mathrm{class}}_{\mathrm{LP}}q
\quad\land\quad
(p\land\neg p)\not\models_{\mathrm{LP}}q.
}
\]
\statuswitness
\end{theorem}

This separation blocks the tempting inference from local anti-diagonal behavior to a
global completeness or consequence-collapse claim. Recovery transfers classical
laws exactly on the recovered domain; it does not remove \(\Bval\) and \(\Nval\)
from the ambient semantics.

\paragraph{Gate boundary.}
Gate~9 proves exact local characterizations of tolerant LEM, strict LEM, and
universal LP explosion; rewrites them through the Gate~8 split representation;
transfers Gate~7 probability recovery to those laws throughout a fixed model; and
identifies a model-relative ST/LP collapse under classical antecedents. An explicit
glut witness separates classical-restricted explosion from unrestricted LP
consequence. No global classical completeness theorem, proof-system equivalence,
arbitrary-CPEL-model theorem, or claim that every formula is recovery-admissible is
made.
